\chapter{Forsøket på ein vitskap om form}

\begin{motto}
Grunnkurset lærte studenten å sjå. Det lærte han aldri korleis\\ blikket kunne grunngjevast, eller overførast utan augo som hadde sett.
\end{motto}

Bauhaus var forsøket på ein vitskap om form. Det vart ein pedagogikk. Førelesinga vart eit grunnkurs, grunnkurset vart ein stil, og stilen reiste verda rundt utan teorien som skulle bere han. Dette er kapitlet der boka si grunntese møter sin sterkaste motstandar, fordi Bauhaus er det næraste designfaget nokon gong har kome eit systematisk, overførbart fundament for forma. Skulen i Weimar, seinare Dessau, til sist Berlin, samla på under femten år dei skarpaste hovuda i europeisk formgjeving og sette dei til å gjere nett det denne boka etterlyser: å finne grunnlaget for forma, ikkje smaken for henne. At dei mislukkast — at dei i staden produserte den mest innflytelsesrike pedagogikken og den mest gjenkjennelege stilen i moderne design — er ikkje eit teikn på at dei var dårlege. Det er det sterkaste provet boka har. Dei beste prøvde, med dei beste føresetnadene, og fundamentet kom likevel ikkje.

\section{Manifestet og ambisjonen}

Då Walter Gropius i 1919 slo saman kunstakademiet og kunsthandverksskulen i Weimar til Das Staatliche Bauhaus, opna han med eit manifest som er verdt å lese som det det er: ei programerklæring om grunnlag. «Det endelege målet for all biletkunst er bygningen,» skriv \textcite{gropius1919}, og kallar arkitektar, målarar og bilethoggarar til å vende seg mot handverket, for «det finst ingen vesensskilnad mellom kunstnaren og handverkaren». Lina, kjend frå Lyonel Feininger sitt tresnitt av ein katedral på framsida, er at alle kunstartane skal sameinast i det store byggverket, slik dei var det i middelalderens bygghyttar. Legg merke til to ting. For det fyrste er dette eit forsøk på å lækje nett den striden Werkbund ikkje fekk avgjort: skiljet mellom kunst og handverk skal opphevast, ikkje ved røysting, men ved ein ny samlande praksis. For det andre er språket framleis Arts and Crafts sitt — handverket, bygghytta, det samlande verket — same kor mykje skulen seinare skulle vende seg mot industrien.

Her ligg alt den indre spenninga som skulle drive heile Bauhaus si historie, og som forklarer kvifor fundamentet aldri sette seg. Skulen var frå starten kløyvd mellom tre ulike svar på kva forma skulle grunnast i: kunsten, handverket og industrien. Den fyrste fasen i Weimar, prega av den karismatiske og mystisk orienterte Johannes Itten, helte mot kunsten og det indre, åndelege. Den seinare fasen, særleg etter flyttinga til Dessau i 1925, helte mot industrien og det rasjonelle, under parolen «kunst og teknikk — ei ny eining». Skulen skifta altså ikkje berre lærarar, men sjølve grunntanken om kva forma var, midt i si korte levetid. \textcite{droste2002bauhaus} dokumenterer desse fasane med ein presisjon som gjer det umogleg å snakke om «Bauhaus» som ein einsarta ting; det var fleire skular under same namn, usamde om det mest grunnleggjande. Eit fag som framleis kranglar om kva det studerer medan det underviser, har ikkje funne grunnen sin.

\section{Vorkurs: pedagogikken som vart kjernen}

Det som overlevde alle fasane, og som vart Bauhaus si verkelege og varige oppfinning, var \emph{Vorkurs} — det obligatoriske grunnkurset alle studentar måtte gjennom før dei fekk velje verkstad. Her, om nokon stad, finn me forsøket på ei systematisk grunnlære om form. Itten utvikla det fyrst, som ei øving i å frigjere studenten frå innlærte konvensjonar og vekkje ei direkte, kroppsleg kjensle for materiale, tekstur, rytme og kontrast. I \textit{Design and Form} legg \textcite{itten1963} fram metoden i ettertid: øvingar i å kjenne motsetnadspar — lyst og mørkt, hardt og mjukt, glatt og ruglete — og å gje dei form. Det er ein pedagogikk av sjeldan kraft, og han verka; studentane kom ut av kurset med eit nytt blikk. Men sjå kva slag kunnskap det var. Itten lærte studenten å \emph{sjå} og å \emph{kjenne}; han lærte han ikkje ein teori om kvifor ein gjeven form var rett, som ein utanforståande kunne etterprøve. Grunnkurset overførte eit blikk, ikkje eit fundament.

Då Itten gjekk i 1923 — etter ein konflikt med Gropius som var like mykje om livssyn som om pedagogikk — vart kurset ført vidare og omforma av László Moholy-Nagy og Josef Albers, og helte no frå det mystiske mot det konstruktive og materielle. I \textit{The New Vision} formulerer \textcite{moholy1947vision} ambisjonen i sin reinaste form: ei ny, objektiv grunnlære om dei visuelle elementa, om forholdet mellom material, rom, lys og rørsle, tenkt nesten som ein grammatikk for det visuelle. Dette er det punktet der Bauhaus kjem aller nærast ein vitskap om form: tanken om at det finst grunnelement og grunnrelasjonar i det visuelle som kan isolerast, lærast og kombinerast etter prinsipp, slik tonar kombinerast i musikk eller ord i ein syntaks. Albers, på si side, dreiv grunnkurset mot det reint materielle: gjev studenten eit ark papir, og la han oppdage gjennom å gjere kva forma kan bere. Det var strålande undervisning. Men spør kva ein sat att med. Eit repertoar av øvingar, ein måte å nærme seg materialet på, eit trena auge — ein \emph{pedagogikk}, overførbar frå lærar til student i rommet, men ikkje ein teori som kunne grunngje ein dom for nokon som ikkje hadde gått kurset.

Skilnaden er fin og avgjerande, så lat oss vere presise om henne. Ein pedagogikk overfører ein dugleik ved øving under rettleiing; ein teori overfører ein kunnskap ved at han er sann og kan prøvast. Når Moholy-Nagy eller Albers sa at ei gjeven løysing var betre enn ei anna, kunne dei vise studenten det — og studenten kom etter kvart til å sjå det same. Men kunne dei \emph{grunngje} det slik at ein som var usamd måtte gje seg? Kunne ein annan lærar, uavhengig, kome til same dom av same grunn? Nei. Dommen kvilte til sjuande og sist på det trena auget, og det trena auget er presis det denne boka frå fyrste kapittel har peika på som ein kompetanse som ikkje let seg skilje frå personen som har han. Bauhaus systematiserte opplæringa av auget meir enn nokon før eller sidan. Men eit systematisk trent auge er framleis eit auge, ikkje eit fundament.

\section{Fargeteorien: tre system, ingen einskap}

Ingen stad er dette tydelegare enn i fargeteorien, der Bauhaus mønstra ikkje éin, men tre store forsøk på ei lære — og der tre stod ved sida av kvarandre utan å bli til éi. Itten utvikla sin fargesirkel og si lære om dei sju fargekontrastane, knytt til ei nesten temperamentslæreaktig tanke om at kvart menneske har sine «eigne» fargar. Wassily Kandinsky, som underviste ved skulen frå 1922, søkte i \textit{Punkt und Linie zu Fläche} ein meir systematisk, nesten lovsøkjande samanheng: at \textcite{kandinsky1926} ville binde bestemte fargar til bestemte former — gult til trekanten, raudt til kvadratet, blått til sirkelen — og bestemte verknader til bestemte element, som om det fanst ein objektiv grammatikk for korleis det visuelle verkar på sjela. Og Josef Albers, seinare, i det som vart hovudverket \textit{Interaction of Color}, gjekk i stikk motsett retning: \textcite{albers1963} viste at fargen ikkje har nokon stabil, objektiv verknad i det heile, men endrar seg fullstendig med kva som ligg attmed — at same fargen ser ut som to, og to fargar ser ut som éin, alt etter konteksten.

Sjå kor djupt desse tre står i strid. Itten grunnar fargen i temperamentet, Kandinsky i ei objektiv kopling mellom form og farge, Albers i den reine relativiteten der ingen farge har ein verknad åleine. Dette er ikkje tre kapittel i éi lære; det er tre uforeinlege teoriar om kva farge \emph{er} og korleis han verkar, framførte av tre lærarar ved same skule, til dels samstundes. Og — dette er poenget — skulen valde aldri mellom dei. Han avgjorde aldri kven som hadde rett, fordi han ikkje hadde nokon prosedyre for å avgjere det. Kandinsky sin påstand om at gult «høyrer til» trekanten er empirisk testbar, og fell ved test; folk koplar ikkje fargar til former på nokon stabil måte. Men Bauhaus testa han ikkje, forkasta han ikkje, og heldt fram med å undervise alle tre systema parallelt. Det er ikkje slik eit fag med eit fundament handsamar motstridande teoriar. Det er slik eit fag utan eit handsamar dei: det legg dei ved sida av kvarandre og kallar mangfaldet ein rikdom. Albers sitt arbeid er forresten det mest varige av dei tre nettopp fordi det er det minst teoretiske og det mest pedagogiske — ei samling demonstrasjonar du må gjere sjølv for å sjå. Igjen: det som overlever frå Bauhaus, er øvinga, ikkje teorien.

\section{Mazdaznan: då fundamentet var ein religion}

Det er ein episode i Bauhaus si tidlege historie som seinare framstillingar gjerne skundar seg forbi, men som boka ikkje har råd til å gå forbi, fordi han syner med brutal tydelegheit kva slags «grunnlag» grunnkurset eigentleg kvilte på i Weimar-åra. Johannes Itten, mannen som forma Vorkurs og dermed sjølve kjernen i Bauhaus-pedagogikken, var ikkje berre lærar; han var tilhengjar — og ved skulen ein slags apostel — for Mazdaznan, ei nyreligiøs lære av zoroastrisk-vegetariansk opphav, importert frå Amerika via Leipzig. Og dette var ikkje ein privatsak han heldt utanfor klasserommet. Mazdaznan styrte Vorkurs innanfrå. Itten påla studentane pusteøvingar og kroppsøvingar før teikninga, for å «reinse» og opne dei; han la om kantina til vegetarkost etter Mazdaznan sine reglar, med kvitlauk og ein særskild graut; somme studentar fasta, raka hovuda og bar ein eigen Itten-designa kuttel. Den åndelege tilstanden til utøvaren vart sett som føresetnaden for den rette forma. Grunnlaget for å sjå rett, slik Itten lærte det, var bokstaveleg talt ei religiøs reinsing av kroppen.

\begin{kasus}{Kvitlauksgrauten i Weimar}
Sjå nøkternt på kva dette inneber. Bauhaus var det fremste forsøket i historia på å gje forma eit systematisk, overførbart grunnlag — og i den fyrste fasen var dette grunnlaget ei pustelære og ein vegetarisk diett henta frå ein amerikansk-zoroastrisk sekt. Det er lett å le av kvitlauksgrauten, men ein bør heller la han stå som eit reint laboratoriebevis. For når ein pedagogikk gjer den \emph{indre tilstanden} til utøvaren til vilkåret for den rette forma — når du må puste rett, ete rett og reinse deg rett for å sjå rett — då har ein sagt høgt og tydeleg at kunnskapen ikkje ligg i forma og ikkje let seg skilje frå personen. Han ligg i ei oppnådd sjelstilstand. Eit slikt grunnlag kan ikkje skrivast ned og sendast til ein framand, for det \emph{er} personen sin tilstand. Itten gjorde berre eksplisitt, og litt for ærleg, det som boka hevdar ligg løynt i heile tradisjonen: at det «trena auget» til sjuande og sist er ein person, ikkje eit prinsipp. At det fremste forsøket på ein vitskap om form opna med ein religion som inngangsport, er ikkje ein pinleg fotnote. Det er diagnosen.
\end{kasus}

Konflikten som tvinga Itten ut i 1923 var difor ikkje først og fremst pedagogisk; han var ein strid om kva slag grunnlag skulen skulle ha. Gropius hadde innsett at ein skule bygd kring ein gurus åndelege reinsing aldri kunne bli den industrielt orienterte institusjonen tida kravde, og at Mazdaznan-regimet gjorde Bauhaus til ein sekt snarare enn ein skule. Då Itten gjekk og László Moholy-Nagy kom inn, var dette eit medvite skifte frå det mystiske til det konstruktive, frå pust til material, frå sjelstilstand til objektiv grunnlære. Frank Whitford skildrar dette skiftet som vendepunktet i Bauhaus si sjølvforståing \autocite{whitford1984bauhaus}, og Éva Forgács viser at det var like mykje ein politisk og ideologisk kamp om kva skulen skulle vere som ein usemje om undervisning \autocite{forgacs1995bauhaus}. Men legg merke til kva som \emph{ikkje} endra seg i overgangen. Moholy-Nagy kvitta seg med kvitlauken og pusteøvingane, men han kunne framleis ikkje grunngje ein gjeven form for nokon utanfor rommet betre enn Itten kunne; han bytte berre ut ein mystikk om sjela med ein retorikk om objektivitet. Auget var framleis eit auge. Det vart berre kledd om frå prest til ingeniør.

\section{Verkstadene og striden faget arva}

Etter grunnkurset gjekk studenten inn i ein verkstad — metall, tre, tekstil, keramikk, veggmåleri, seinare også typografi og reklame. Det var her Bauhaus skulle gjere det Werkbund ikkje fekk til: foreine den kunstnarlege og den industrielle forma i ein praksis, ikkje i ein debatt. Den såkalla doble leiinga, der kvar verkstad fyrst hadde både ein «formmeister» (kunstnaren) og ein «handverksmeister», var sjølve den institusjonelle innrømminga av at skulen ikkje hadde løyst spørsmålet: ein måtte ha begge til stades, fordi ingen visste korleis det eine kunne bli det andre. \textcite{whitford1984bauhaus} skildrar dagleglivet i desse verkstadene med ein varme som ikkje skjuler spenningane; striden mellom kunstnarane som ville skape og industrien som ville produsere, gjekk gjennom kvart rom. Det store kjeldeverket til \textcite{wingler1969bauhaus} samlar dokumenta som viser at skulen sjølv var fullt klar over denne uløyste spenninga, og diskuterte henne uavlateleg utan å nå nokon konklusjon.

Resultatet var i praksis det same valet, gjort om att av kvar verkstad for seg, som Werkbund hadde røysta over: nokre arbeid helte mot det kunstnarlege unikumet, andre mot den industrielle prototypen. Nokre Bauhaus-produkt vart faktisk sette i produksjon — Marianne Brandt sine lampar, Marcel Breuer sine stålrøyrstolar — og desse vert framheva i ettertid som beviset på at syntesen lukkast. Men sjå kor få dei var, og kor mykje av produksjonen som forblei prototypar, kunstverk, eksperiment. Den katalogen \textcite{bergdoll2009bauhaus} sette saman ved Museum of Modern Art, gjev det grundigaste oversynet me har over kva skulen faktisk produserte, og biletet som teiknar seg er ikkje ein triumferande syntese, men ein vedvarande pendling mellom polane. Skulen løyste ikkje striden mellom kunst og industri; han held han levande, og sende han vidare til alle skular som sidan har modellert seg på han. Den striden lever framleis i kvart designstudium i dag, mellom dei som vil at studenten skal uttrykkje seg og dei som vil at han skal løyse eit problem. Bauhaus arva han frå Werkbund og gav han vidare uløyst. Det er ikkje overføring av eit fundament; det er overføring av eit ope spørsmål.

\section{«Kunst und Technik» og rekneskapen som ikkje gjekk opp}

Den offisielle forteljinga om Dessau-Bauhaus er at skulen der løyste det Weimar ikkje fekk til: under parolen «Kunst und Technik — eine neue Einheit», kunngjord av Gropius i 1923, vende skulen seg mot industrien, og verkstadene byrja for alvor å produsere prototypar for serietilverking. Det er den triumferande versjonen, og ho byggjer på nokre få, mykje reproduserte ikon — Marcel Breuer sin stålrøyrstol, Marianne Brandt og Wilhelm Wagenfeld sine lampar, Wagenfeld sitt glasservise. Men granskar ein lisensieringsrekneskapen til skulen, sprekk forteljinga. Svært få av desse prototypane gjekk faktisk i industriell produksjon i Bauhaus si levetid, og endå færre tente pengar. Verkstadene var, gjennom det meste av historia, tapsprosjekt som måtte subsidierast; den reine, sakleg-industrielle Bauhaus-forma var dyr å lage og selde dårleg. Magdalena Droste, som har gått gjennom skulen sitt eige materiale grundigare enn dei fleste, dokumenterer kor stort gapet var mellom den industrielle retorikken og den faktiske produksjonen: «Kunst und Technik» var ein ambisjon og ein parole, ikkje ein realisert praksis \autocite{droste2002bauhaus}.

Og her kjem ironien som er for god til å gå forbi, fordi han seier alt. Det \emph{eine} Bauhaus-produktet som vart ein verkeleg kommersiell suksess i serie, var ikkje ein stålrøyrstol eller ei ikonisk lampe. Det var tapet. I 1929 inngjekk skulen ein avtale med fabrikanten Emil Rasch om ein serie Bauhaus-tapet — diskré, einsfarga, lett strukturerte flater, utan mønster i det heile. Dei selde i millionar av rullar, vart den langt største inntektskjelda skulen nokon gong hadde, og overlevde skulen sjølv med tiår. Tenk på kva dette tyder for ein institusjon som opna med Gropius sitt manifest om at all kunst kulminerer i det store byggverket, og som sette dei skarpaste hovuda i Europa til å finne grunnlaget for forma. Den varige, lønsame, industrielle suksessen deira var ein vegg utan motiv. Det er ikkje ein hån mot Bauhaus; det er ein presis figur. Det skulen kunne masseprodusere med hell, var nett det som kravde minst av ein teori om form — ei tom, jamn flate — medan dei ambisiøse, teoritunge gjenstandane forblei prototypar og museumsstykke.

Barry Bergdoll og Leah Dickerman sin store katalog frå Museum of Modern Art stadfestar dette biletet når ein les han mot myten: den faktiske produksjonsrekorden er ein vedvarande pendling og eit knippe ikon, ikkje ein industriell revolusjon \autocite{bergdoll2009bauhaus}. Poenget for boka er ikkje at Bauhaus mislukkast kommersielt — mange gode ting tener lite. Poenget er at sjølv \emph{suksessen} stadfestar tesen. Når eit fag som har lova ein vitskap om form endar med at det reproduserbare produktet er ein motivlaus tapet, medan det formteoretisk ambisiøse blir verande unika, då har marknaden uforvarande felt den same dommen som boka feller: det fanst ein pedagogikk, det fanst ein stil, og det fanst eit auge — men det fanst ingen teori om form som kunne gjere det vanskelege reproduserbart. Tapeten kunne reproduserast fordi han ikkje trong teorien. Stolen vart eit ikon, ikkje ei serievare, fordi han trong meir grunngjeving enn skulen hadde å gje.

\section{Politikken, stenginga, og det som reiste vidare}

At Bauhaus aldri fekk fullføre nokon syntese, skuldast ikkje berre indre uvisse, men ytre vald. Skulen var politisk omstridd frå fyrste dag — for venstrevridd, for internasjonal, for «kulturbolsjevikisk» i auga til den tyske høgresida. Han vart pressa ut av Weimar i 1925, fann tilflukt i Dessau, vart pressa ut av Dessau i 1932 då nasjonalsosialistane fekk makta i delstaten, opna eit siste, kort kapittel som privat skule i Berlin under Ludwig Mies van der Rohe, og stengde seg sjølv i 1933 under press frå Gestapo. \textcite{forgacs1995bauhaus} viser at politikken ikkje var ein ytre ramme rundt skulen, men gjennomsyra sjølve ideen og kvardagen hans; striden om kva Bauhaus skulle vere var alltid òg ein strid om kva slags samfunn forma skulle tene. Det er ein viktig korreksjon til den reinsa, apolitiske Bauhaus-myten som seinare vart eksportert til Amerika.

For det var eksporten som avgjorde ettermælet. Då lærarane spreidde seg i emigrasjon — Gropius og Breuer til Harvard, Mies og Moholy-Nagy til Chicago, Albers til Black Mountain og sidan Yale — tok dei med seg det av Bauhaus som let seg ta med. Og legg merke til kva det var. Det var ikkje teorien, for han fanst ikkje som ein avslutta teori. Det var grunnkurset — Vorkurs vart \emph{the foundation course} ved kunstskular over heile den vestlege verda — og det var stilen, det reine, geometriske, sakleg-elegante Bauhaus-uttrykket, som vart sjølve biletet på det moderne og spreidde seg til møblar, bygningar, bøker og skrifttypar i kvar verdsdel. Pedagogikken reiste, og stilen reiste. Teorien, som skulle bere dei begge, stod att i Weimar og Dessau som eit uavslutta forsøk. Dette er det presise innhaldet i påstanden som opna kapitlet: førelesinga vart eit grunnkurs, grunnkurset vart ein stil, og stilen reiste verda rundt utan teorien som skulle bere han.

Det er lærerikt å sjå sidelengs, mot aust, på det som skjedde samstundes. I Sovjet-Russland arbeidde Vchutemas-skulen og konstruktivistane med ein på mange måtar parallell ambisjon: ei objektiv, materialistisk grunnlære om form, tufta på vitskap og produksjon snarare enn på kunstnaren sitt indre. \textcite{lodder1983} viser at konstruktivismen, i si strengaste form, faktisk dreiv ambisjonen om ein vitskap om form lenger enn Bauhaus gjorde — mot ein reell teori om «fakturen», om materialet og konstruksjonen sin eigen lovmessigheit. Og likevel enda det same: ein revolusjonær estetikk, ein pedagogikk, eit knippe meisterverk, men ikkje eit overført, kumulativt fundament, og dessutan ein brutal politisk slutt då Stalin kvelte avantgarden. To av dei mest seriøse forsøka i historia på å gjere form til vitskap, eitt i vest og eitt i aust, leverte begge ein stil og ein pedagogikk, og ingen av dei eit fundament. At det skjedde to gonger, uavhengig, med dei beste kreftene, er ikkje eit samantreff. Det er strukturen i saka.

\textcite{banham1960} sette til slutt Bauhaus inn i den breiare soga om teori og form i den fyrste maskinalderen, og dommen hans er edrueleg: det modernismen kalla teori, var oftast ein rasjonalisering i etterhand av val gjorde på estetisk grunnlag. Bauhaus gav desse vala den mest gjennomarbeidde pedagogiske forma dei nokon gong fekk. Men ein pedagogikk som lærer deg å gjere dei rette vala, utan å kunne grunngje kvifor dei er rette for nokon utanfor rommet, er ikkje ein vitskap. Han er ein lærdomstradisjon — den fremste designfaget har hatt, og likevel berre ein lærdomstradisjon. Det Bauhaus etterlét, var ein måte å sjå på, og ein måte å undervise det blikket på. Det var mykje. Det var ikkje grunnen.

\vignett

Bauhaus var det høgaste punktet for ambisjonen om ein vitskap om form, og difor den klaraste målinga av kor langt ambisjonen nådde: til ein pedagogikk og ein stil, og ikkje eit steg lenger. Men sjølve forsøket på å gjere form til vitskap var ikkje dødt med skulen. Det skulle reise seg att, reinsa for kunstnaren og det indre, ved ein ny skule i ein liten tysk by tjue år seinare — ein skule som ville ta Bauhaus si uavslutta oppgåve og endeleg gjere ho ferdig med metode, system og vitskap. Det heitte Ulm, og det er dit historia no går.

\laeringsmal{\begin{itemize}
\item Forklare kvifor Bauhaus er det sterkaste provet for boka sin tese: dei beste prøvde, med dei beste føresetnadene, og fundamentet kom likevel ikkje.
\item Skilje mellom ein \emph{pedagogikk} (overfører ein dugleik ved øving) og ein \emph{teori} (overfører ein kunnskap ved at han kan prøvast), og bruke skiljet på Vorkurs.
\item Bruke Mazdaznan-episoden og dei tre uforeinlege fargeteoriane til å vise at grunnlaget kvilte på det trena auget, ikkje på eit prinsipp.
\item Gjere greie for gapet mellom «Kunst und Technik»-retorikken og den faktiske produksjonsrekorden, og kva Rasch-tapeten seier om kva som let seg reprodusere utan ein teori.
\end{itemize}}

\ovingar{\begin{enumerate}
\item Bauhaus underviste tre uforeinlege fargeteoriar (Itten, Kandinsky, Albers) samstundes utan å velje mellom dei. Korleis handsamar eit fag med eit fundament motstridande teoriar? Finn eit døme frå eit anna fag og samanlikn.
\item Itten gjorde den indre tilstanden til utøvaren til vilkår for den rette forma. Finn ein samtidig designpedagogikk eller «kreativitetsmetode» som gjer noko liknande (mindfulness, flow, «design mindset»). Er strukturen den same?
\item Det eine Bauhaus-produktet som vart ein serieсуksess var ein motivlaus tapet. Diskuter kvifor nett det som kravde minst teori, lét seg reprodusere best. Gjeld den same logikken i dag?
\item Vorkurs vart «the foundation course» ved kunstskular over heile verda, men teorien som skulle bere han følgde ikkje med. Vurder eit grunnkurs du sjølv har teke: overførte det eit blikk eller eit fundament? Korleis veit du skilnaden?
\end{enumerate}}

\vidarelesing{\fullcite{wingler1969bauhaus}; \fullcite{whitford1984bauhaus}; \fullcite{droste2002bauhaus}; \fullcite{bergdoll2009bauhaus}; \fullcite{forgacs1995bauhaus}; \fullcite{lodder1983}; \fullcite{banham1960}.}
